(2 points + 4 points)
\begin{teilaufgaben}
\item
The function $u(x,t)$ is defined on the upper half plane
$\Omega = (-\infty, \infty) \times [0,\infty)$.

The function $u(x,t)$ satisfies in $\Omega$ the equation
\[
u_{x}(x,t) + u_{t}(x,t) = 0
\]
and $u(x,0) = x^3+17$ on the boundary of $\Omega$. 

Determine an approximate value for $u(0,1)$.

Use the scheme of Lax-Wendroff with $\Delta x = 0.1$ and $\Delta t = 0.1.$

\item
The function $u(x,t)$ is defined on the strip
$\Omega = [0, 3] \times [0,\infty)$.

The function $u(x,t)$ satisfies in $\Omega$ the equation
\[
u_{xx}(x,t) = u_{tt}(x,t)
\]
and  $u(0,t) = 0$, $u(3,t) = t$ on the boundary of $\Omega$. 

The function $u(x,t)$ also satisfies the initial conditions
\[
u(x,0)
=
0
\qquad \text{and} \qquad
u_t(x,0)
=
1/3 \cdot (x^3 - 6 \cdot x^2 + 12 \cdot x - 6).
\]

Determine approximate values for $u(1,1)$ and $u(2,1)$  

Use finite differences with $\Delta x = 1$ and $\Delta t = 1/2.$
\end{teilaufgaben}

\begin{loesung}
\begin{teilaufgaben}
\item
Set
\[
U_{k,j} = \tilde u(k/10, j/10).
\]
Apply the Lax-Wendroff scheme with $r = 1$ to obtain:
\begin{equation}
U_{k+1,j+1} = U_{k,j}.
\tag{\textbf{1P}}
\end{equation}
Therefore:
\begin{equation}
\tilde u(0,1) = 17 + (-1)^3 = 16.
\tag{\textbf{1P}}
\end{equation}

\item
We have $f(x) = 0$ and $g(x) = 1/3 \cdot (x^3 - 6 \cdot x^2 + 12 \cdot x - 6)$.
Hence for the first time step we have by Taylor
\begin{equation}
\tilde u(1, 1/2)
=
1/3 \cdot 1 \cdot 1/2 = 1/6
\qquad
\text{and}
\qquad
\tilde u(2, 1/2) = 1/3 \cdot 2 \cdot 1/2 = 1/3.
\tag{\textbf{1P}}
\end{equation}

For the second time step we use the leapfrog scheme with $r = 1/2$. We have
 
\[
\tilde u(1, 1)
=
1/4 \cdot \tilde u(0, 1/2) + 6/4 \cdot \tilde u(1,1/2) + 1/4
\cdot 
\tilde u(2, 1/2) - \tilde u(1,0) 
\]
and
\begin{equation}
\tilde u(2, 1)
=
1/4 \cdot \tilde u(1, 1/2)
+
6/4 \cdot \tilde u(2,1/2)
+
1/4 \cdot  \tilde u(3, 1/2)
-
\tilde u(2,0).
\tag{\textbf{2P}}
\end{equation}
This leads to
\begin{equation}
\tilde u(1, 1) = 1/3, \quad \tilde u(2, 1) = 2/3
\tag{\textbf{1P}}
\end{equation}

\end{teilaufgaben}
\end{loesung}


